\firstword{T}{he} needs for encryption quickly emerged in human's history, mostly because of military operations or to exchange sensitive informations. However, nowadays, with the development of Internet, everybody needs to exchange encrypted data : texts, videos, images.. 
Most of the traditional, and accepted, symmetric algorithms (AES-128, RC6..) to encrypt data are block-cipher algorithm. This type of algorithm lead to large computational time to encrypt large images. Moreover, the redundancy of information in image data lead some encryption algorithm to show some weaknesses when it comes to image encryption.  

Chaos-based cryptography can be suitable to encrypt images for different reasons : with a chaotic map, a key can be generated in no time, which lead to reasonable computational time. The chaotic map being sensible to tiny variations in initial conditions lead the key sensible to those variations, and consequently, the encryption. It remains, however, a question, which is fundamental in cryptography : is it safe to encrypt images based on chaos theory ? 

In this paper, we will present an algorithm for image encryption, and evaluate it with basic metrics as visual, histogram and correlation analysis. \\
The paper is organized in five sections : the \textit{Section 2} introduces the reader to the algorithm used to encrypt images. A basic analysis of the properties in regards of security is given in \textit{Section 3}. We finally conclude on this paper in \textit{Section 4}. \textit{Section 5} is reserved to mathematical details.

